\chapter{Deux grands projets découverts}
\label{chap:projets}

Pendant mon stage, deux projets du Groupe ont particulièrement retenu mon attention~: la mine d'or de Boto, au Sénégal, et le projet cuprifère de Tizert, au Maroc, sur lequel a porté le travail de suivi auquel j'ai été associée. Les données chiffrées de ce chapitre proviennent de sources publiques et sont référencées~; ce que j'ai personnellement observé est signalé comme tel.

% ---------------------------------------------------------------------
\section{Boto --- une mine d'or au Sénégal}

\subsection{Localisation et objectifs}

Le projet Boto est situé dans la région de Kédougou, à l'est du Sénégal, près des frontières du Mali et de la Guinée, dans une zone géologique (le \textit{Senegalo-Malian Shear Zone}) qui abrite plusieurs gisements aurifères de classe mondiale~\cite{managem_boto}. Il s'agit d'une mine d'\textbf{or} à ciel ouvert, acquise par Managem en 2023 auprès de la société canadienne IAMGOLD, le Groupe en détenant 90\,\% aux côtés de l'État sénégalais~\cite{allafrica_boto}.

\subsection{Un investissement majeur et un statut récent}

Boto représente un investissement de l'ordre de 350~millions d'euros, ce qui en fait l'un des plus importants développements miniers du Sénégal~\cite{managem_boto,northafrica_boto}. Le projet vise une production moyenne d'environ 160\,000~onces d'or par an sur ses trois premières années, pour des réserves estimées à près de 1,8~million d'onces et une durée de vie d'une douzaine d'années~\cite{managem_boto}. La mine s'appuie sur une usine de traitement et des infrastructures dédiées (centrale électrique, digues, base-vie, piste d'atterrissage, réseau routier)~\cite{northafrica_boto}.

Le fait le plus marquant est que Boto est passé, très récemment, du stade de la construction à celui de la \textbf{production}~: la première coulée d'or a eu lieu le 2~septembre 2025~\cite{managem_boto}. Le projet est donc, à l'échelle du Groupe, un exemple actuel de passage réussi « du chantier à l'exploitation ».

\subsection{Importance et défis d'ingénierie}

Pour Managem, Boto consolide la stratégie aurifère du Groupe en Afrique de l'Ouest et renforce son positionnement de producteur régional d'or~\cite{allafrica_boto}. Sur le plan de l'ingénierie, un tel projet illustre bien les enjeux évoqués par les ingénieurs~: construire, dans une zone isolée, non seulement une usine mais tout un écosystème (énergie, eau, logistique, hébergement), et le faire dans des délais maîtrisés. Je n'ai pas travaillé sur ce projet~; il m'a servi de \emph{point de repère} pour comprendre ce qu'est un projet minier arrivé à maturité.

% ---------------------------------------------------------------------
\section{Tizert --- le projet cuprifère}

\subsection{Un projet stratégique pour le cuivre marocain}

Le projet Tizert est situé dans l'Anti-Atlas occidental, à environ 65~km de Taroudant, sur la bordure nord de la boutonnière d'Igherm~\cite{managem_tizert}. C'est l'un des plus importants gisements de \textbf{cuivre} (Cu-Ag) du Maroc, exploité en fosse et en souterrain. Le projet représente un investissement de l'ordre de 440~millions de dollars et vise une capacité de production de l'ordre de 100 à 120~kilotonnes de concentré de cuivre par an, pour une durée de vie d'environ 17 à 18~ans~\cite{managem_tizert}. Comme Boto, Tizert est entré en production commerciale en 2025, avec les premières ventes de concentré au quatrième trimestre~\cite{managem_tizert}.

L'importance de Tizert dépasse le seul cadre du Groupe~: sa production doit couvrir une part significative des besoins nationaux en cuivre, métal essentiel à la transition énergétique~\cite{managem_tizert}. Les études de faisabilité du projet ont d'ailleurs mobilisé Reminex aux côtés de cabinets d'ingénierie de renommée internationale~\cite{managem_tizert}, ce qui donne un sens très concret au rôle de la filiale décrit précédemment.

\subsection{L'installation des convoyeurs ROM}

C'est sur un chantier précis de Tizert que j'ai été associée à un travail de suivi~: l'\textbf{installation des convoyeurs ROM} (\textit{Run-Of-Mine}), c'est-à-dire les convoyeurs à bande qui acheminent le minerai brut vers l'unité de traitement. D'après les documents de suivi que j'ai contribué à mettre en forme, ce chantier porte sur deux lignes de convoyage identifiées par leur numéro de plan~:

\begin{itemize}
\item le \textbf{convoyeur principal (Plan~1005)}, une ligne aérienne suspendue (tronçons, tour d'escalier, structure treillis, pylônes de support)~;
\item le \textbf{convoyeur du tunnel (Plan~1037)}, une ligne en galerie.
\end{itemize}

L'avancement est suivi selon trois jalons cumulatifs --- inventaire, pré-assemblage et montage --- et fait l'objet d'un tableau de bord hebdomadaire (indicateurs de sécurité, tonnage réalisé, courbes d'avancement, planning). Un exemple de ce rapport, dont j'ai contribué à la mise en forme, est reproduit en annexe~\ref{ann:rapport} \emph{à titre d'illustration du format de reporting}, et non comme un travail dont je serais l'auteur.

\begin{presente}
\textbf{Mon rôle, précisément.} Je n'ai pas conçu ni dimensionné ces convoyeurs. J'ai aidé à \emph{préparer} les rapports d'avancement~: collecter et organiser les informations transmises par les équipes de montage, mettre en forme les tableaux et les figures, et produire un document professionnel. Ce travail, réalisé sous la supervision de mon tuteur, est détaillé au chapitre~\ref{chap:stage}.
\end{presente}
